\section{Related work}
\label{sec:related}

\paragraph{Agent-memory benchmarks.} LoCoMo \citep{maharana2024locomo}, LongMemEval \citep{wu2025longmemeval},
MemoryAgentBench \citep{hu2025memoryagentbench} and PrefEval \citep{zhao2025prefeval} are the standard yardsticks for
memory systems such as MemGPT \citep{packer2023memgpt}, Mem0 \citep{chhikara2025mem0}, Zep \citep{rasmussen2025zep},
A-MEM \citep{xu2025amem} and Nemori \citep{nemori2025}. Newer benchmarks avoid the prior by construction: they use
fictional worlds \citep{groundtruthfirst2026,memgym2026}, counterfactual edits that contradict it
\citep{hu2025memoryagentbench}, or drop items a model answers without context
\citep{longmemevalvtwo2026,recon2026}. \citet{groundtruthfirst2026} state the principle that a memory score partly
measures the model's priors, and \citet{memgym2026} show that without fictionalisation frontier models score highly on their benchmark from
pretraining alone. PersonaMem-v2 \citep{personamemv2} reports that personalisation accuracy falls from stereotypical to
neutral to anti-stereotypical preferences and attributes this to population priors; \citet{feng2026knowit} find that
this gap, present for a bare backbone, reverses once memory is retrieved. On the write side, Nemori's Predict-Calibrate
step stores what the model failed to predict, an explicitly prior-aware write policy; in our controlled study it kept
every item (\S\ref{sec:controlled}). We differ in three ways: we audit and re-score the benchmarks already in use, we give a
scoring procedure with explicit assumptions and diagnostics rather than a filter, and we use same-slot minimal pairs that
hold the item fixed while manipulating its congruence with the prior. The handling of LoCoMo's unanswerable category is
known to be inconsistent across papers and to move headline scores \citep{acm2026}; we show that its
accuracy is inversely related to how much a system remembers.

\paragraph{Locating memory failures.} Stage-wise analyses attribute errors to parts of the memory pipeline: extraction
and updating \citep{chen2025halumem}, indexing, retrieval and reading \citep{wu2025longmemeval,yuan2026diagnosing}, or
knowing versus acting on a stated preference \citep{feng2026knowit}. They find that systems which extract before storing
lose much of what they are told (recall below 60\% for every system tested by \citealp{chen2025halumem}), that retrieval
dominates once raw history is kept \citep{yuan2026diagnosing}, and that 18\% of one system's LoCoMo errors are abstentions
despite retrieved evidence \citep{shu2026remem}, the suppression that leads us to compute \cpr{} under a forced protocol.
Error taxonomies built from wrong answers \citep{yuan2026diagnosing,shu2026remem} never see a stage failure that a correct
guess rescues; LongMemEval, which checks retrieval for every question, reports that correct retrieval is necessary for a
correct answer ``in $\sim$90\% of the time''. \pcme{} is complementary: it needs only final answers and the reader's guess
rate, it does not say which stage failed, and it identifies the correct answers that carry no evidence of memory.

\paragraph{Evidence-free baselines and prior knowledge.} Question-only and hypothesis-only baselines revealed that
reading-comprehension and NLI benchmarks could be solved without the passage or the premise
\citep{kaushik2018much,poliak2018hypothesis,gururangan2018annotation}; a no-context baseline plays the same role for
long-context QA \citep{longbench100}. Work on knowledge conflicts studies when models follow a counterfactual context
over their parametric knowledge \citep{longpre2021entity,xie2024adaptive,wu2024clasheval,ming2025faitheval}. Our counter-schema items are a memory-side
analogue: the fact is stated once, possibly many sessions before it is needed, and must survive the memory before it can
override the prior.

\paragraph{Guessing corrections and serial reproduction.} Correcting test scores for guessing has a long history in
psychometrics \citep{lord1968statistical}; our guess-corrected score generalises formula scoring to item-specific guess
rates estimated from the model itself. Our controlled study borrows the logic of serial reproduction
\citep{bartlett1932remembering} and of iterated learning, where transmission chains converge to the learners' prior
\citep{griffiths2007language,kirby2007innateness}. We find that strong compactors do not rewrite toward the prior; the
prior re-enters when the reader fills the gaps that deletion leaves, a memory-side counterpart of the unsupported user
attributes that personalised models add on their own \citep{sun2026mirage}.
